% generated by tools/build_cv.py from assets/data.js: edit the data, not this file
\resumeSubHeadingListStart
  \resumeSubheading{Probing Accretion Disk Sizes with Microlensing of Lensed AGN Light Curves from LSST}{Apr. 2025 -- Present}{\textit{Collaborators: \cvlink{https://sibirrer.github.io}{Prof. Simon Birrer}, Dr. Joshua Fagin, Dr. Henry Best, and others}}{}
  \resumeItemListStart
    \resumeItem{Using simulations to train a simulation-based-inference forward model that learns how microlensing alters the multi-band light curves of strongly lensed AGN, with realistic LSST cadences, time-variable microlensing, and realistic source-quasar and deflector populations.}
    \resumeItem{Aiming to constrain the temperature power-law slope, black-hole mass and wavelength dependence of disk size for AGN accretion disks in the thin-disk model.}
  \resumeItemListEnd
  \resumeSubheading{Constraining Dark Energy with Pseudo-Double-Source-Plane Strong Gravitational Lenses}{Apr. 2025 -- Present}{\textit{Supervisor: \cvlink{https://sibirrer.github.io}{Prof. Simon Birrer}, Department of Physics \& Astronomy, Stony Brook University}}{}
  \resumeItemListStart
    \resumeItem{Pairing single-source-plane galaxy--galaxy lenses with near-identical deflectors to mimic double-source-plane lenses, which avoids the intermediate source's mass and its mass-sheet degeneracy.}
    \resumeItem{With the LSST Y10 galaxy--galaxy lens sample, forecast constraints on flat w$_0$w$_a$CDM reach $\sigma$(w$_0$) $\approx$ 0.27.}
    \resumeItem{Establishes pseudo-DSPLs and strong lensing tomography as a complementary probe of dark energy.}
  \resumeItemListEnd
  \resumeSubheading{SLSim: Strong Lensing Simulation Pipeline}{}{\textit{Co-lead developer and maintainer, with the LSST Strong Lensing and Dark Energy Science Collaborations}}{}
  \resumeItemListStart
    \resumeItem{Open-source package that simulates populations of strong lenses (galaxies, lensed quasars, lensed supernovae) as Rubin LSST, Euclid, Roman and other observatories will see them. \cvlink{https://github.com/LSST-strong-lensing/slsim}{github.com/LSST-strong-lensing/slsim}}
  \resumeItemListEnd
  \resumeSubheading{Multi-Band Linear Polarimetry of the Solar Corona during Total Solar Eclipses}{Apr. 2023 -- 2026}{\textit{Supervisor: \cvlink{https://supritsingh.github.io}{Prof. Suprit Singh}, Department of Physics, IIT Delhi}}{}
  \resumeItemListStart
    \resumeItem{Observed the total solar eclipses of 20 April 2023 (Exmouth, Australia) and 8 April 2024 (Dallas, Texas), imaging the corona in red, green and blue through polarizers at $-$60$^\circ$, 0$^\circ$ and +60$^\circ$ over a nine-step exposure bracket.}
    \resumeItem{Reduction pipeline: registration by lunar-limb fitting and plate solving, photometrically linear HDR synthesis, and multi-scale Gaussian normalization with polarization mapped to colour.}
  \resumeItemListEnd
  \resumeSubheading{Understanding Galactic Dynamics with Pattern Speed Estimation}{Aug. 2023 -- July 2024}{\textit{Supervisor: \cvlink{https://www.astro.sunysb.edu/koda/Stony_Brook/Jin_KODA.html}{Prof. Jin Koda}, Department of Physics \& Astronomy, Stony Brook University}}{}
  \resumeItemListStart
    \resumeItem{Analysed the variation of the pattern speed in N-body hydrodynamic simulations of a Milky-Way-like galaxy.}
    \resumeItem{Tested pattern-speed estimators such as the Tremaine--Weinberg method, and modifications of it, for real observations of barred spiral galaxies.}
  \resumeItemListEnd
  \resumeSubheading{Estimating Galactic Magnetic Fields using Mg II Absorbers and Background Quasars}{Aug. 2022 -- Aug. 2024}{\textit{Supervisors: \cvlink{https://supritsingh.github.io}{Prof. Suprit Singh} (IIT Delhi) and Dr. Sunil Malik (Universität Potsdam)}}{}
  \resumeItemListStart
    \resumeItem{Matched the NVSS rotation-measure catalogue of 37K sources to 750K quasars from SDSS DR16.}
    \resumeItem{Probed line-of-sight magnetic fields in galaxies at 0.37 $\leq$ z $\leq$ 0.8 through the excess residual rotation measure of quasar sightlines with and without intervening Mg II absorbers.}
  \resumeItemListEnd
\resumeSubHeadingListEnd
